\section{Introduction}
\label{sec:tpc34-introduction}

The zero-frequency coefficient left by the matched cutoff construction is
no longer an unspecified parity remainder.  TPC-32 identifies the exact
distinguished coefficient and the \(1/400\) endpoint ledger
\citep{WangTPC32}; TPC-33 returns to the physical row mask, opens a
primitive affine column, proves a fixed-determinant finite-field second
moment, and reduces the scalar shell to the energy condition
\begin{equation}
 E_L+E_R\ll_{\eps,h,\mathscr D}X^\eps Q^3.
 \label{eq:tpc34-inherited-Q3-gate}
\end{equation}
The implication from \eqref{eq:tpc34-inherited-Q3-gate} to the required
\(Q^2\) scalar amplitude is proved in TPC-33 and is recalled in
\cref{sec:tpc34-physical-interface}.  The estimate itself is not inherited
\citep{WangTPC33}.

The purpose of this paper is to determine what a literal \(TT^*\)
analysis does and does not prove at that gate.  This is useful for two
reasons.  First, the column energy hides two independent duplications: a
row duplication and an orbit-time duplication.  Treating both as one
generic off-diagonal discards a small but decisive endpoint margin.
Second, opening the actual row-gcd projector creates genuine divisor
incidence, but the principal divisor \(v=1\) survives with no gain.  The
remaining obstruction can therefore be stated as a concrete averaged
M\"obius row energy rather than as a vague request for more
pseudorandomness.

The algebraic \(TT^*\) principle and the large sieve are classical
\citep{Bombieri1965,IwaniecKowalski2004}.  The contribution here is not a
new abstract Hilbert-space inequality.  It is the exact placement of that
inequality inside the physical matched shell, including the true mask,
the ultra-divisor variable, the endpoint exponents, and the sign-sensitive
terminal layer.

\subsection{Main results}

Let
\[
 K^L_{\gamma,\alpha}
 =\sum_j\mathfrak A^{\rm act}_{\alpha,\gamma}(j)
 C_{m_\alpha}(j)\mathsf H_{m_\gamma}(j),
 \qquad
 E_L=\|K^L\gamma^{(1)}\|_2^2,
\]
with \(K^R,E_R\) defined by exchanging the two rows and the two cutoff
legs.  Opening both the ultra increment \(C\) and the row-gcd projector
gives an exact synthesis matrix \(T_L=(\Phi_\gamma(\xi))\), where
\(\xi=(\alpha,j,u,v)\), and
\begin{equation}
 E_L
 =\sum_{\xi_1,\xi_2}
 \overline{z_L(\xi_1)}z_L(\xi_2)
 \cG_L(\xi_1,\xi_2),
 \qquad \cG_L=T_L^*T_L.
 \label{eq:tpc34-intro-Gram}
\end{equation}
The physical opposite-row count gives
\begin{equation}
 |\cG_L(\xi_1,\xi_2)|
 \ll X^\eps\left(L+\frac{Q}{[v_1,v_2]}\right).
 \label{eq:tpc34-intro-kernel-bound}
\end{equation}
This is strongest when the two gcd-projector divisors have a large least
common multiple and is trivial on \(v_1=v_2=1\).

The first positive theorem closes the entire Gram-copy
(same-input-row) diagonal:
\begin{equation}
 E_{L,\Delta}+E_{R,\Delta}
 \ll X^\eps Q^2J^2.
 \label{eq:tpc34-intro-row-diagonal}
\end{equation}
When \(Q=X^\beta\), \(J=X^{1-\beta}\), this crosses the \(Q^3\) gate as
soon as \(\beta>2/3\).  At the physical endpoint
\(\beta=267/400\), its ratio to \(Q^3\) is exactly
\begin{equation}
 \frac{Q^2J^2}{Q^3}=\frac{J^2}{Q}
 =X^{-1/400+o(1)}.
 \label{eq:tpc34-intro-margin}
\end{equation}
The margin is small, but it is a real power and no conjectural input is
used.

Two further collision estimates locate the remaining mass.  The absolute
Gram mass with equal targets,
\(m_{\alpha_1}j_1=m_{\alpha_2}j_2\), is
\begin{equation}
 \ll X^\eps Q^2J.
 \label{eq:tpc34-intro-target-collision}
\end{equation}
The absolute mass with the same ultra divisor \(u_1=u_2>T\gg J\) is
\begin{equation}
 \ll X^\eps\frac{Q^3J^2}{T}.
 \label{eq:tpc34-intro-ultra-collision}
\end{equation}
Thus shared-divisor incidence saves the factor \(T\) from the absolute
baseline but still misses the \(Q^3\) gate by \(J^2/T\).  On the same
row, however, a shared ultra divisor forces the orbit times to coincide,
because
\begin{equation}
 (mj+h,mj'+h)=(mj+h,j-j').
 \label{eq:tpc34-intro-transposed-gcd}
\end{equation}

The cleanest next interface is obtained before duplicating orbit time.
Define \(V_L,V_R\) by freezing \(j\) inside the row sum.  Then
\begin{equation}
 E_L+E_R\ll X^\eps J(V_L+V_R).
 \label{eq:tpc34-intro-orbit-transfer}
\end{equation}
Hence
\begin{equation}
 V_L+V_R\ll X^\eps Q^3/J
 \label{eq:tpc34-intro-V-gate}
\end{equation}
is sufficient for \eqref{eq:tpc34-inherited-Q3-gate}.  The Gram-copy diagonal of
this new energy is \(O(X^\eps Q^2J)\), and
\begin{equation}
 \frac{Q^3/J}{Q^2J}=\frac{Q}{J^2}
 =X^{1/400+o(1)}.
 \label{eq:tpc34-intro-V-allowance}
\end{equation}
The permissible total is therefore only a small power above the natural
diagonal scale.  Proving the stronger diagonal-scale estimate
\(V_L+V_R\ll X^\eps Q^2J\) would return
\[
 E_L+E_R
 \ll X^\eps Q^2J^2
 \ll_{\mathscr D}X^{2+\eps}
 =X^{\eps+o(1)}Q^3X^{-1/400}.
\]

\subsection{The terminal sign boundary}

On the high--\(\beta\) endpoint packet,
\eqref{eq:tpc34-terminal-range-condition} gives
\(T<N_m(j)<U_0\).  Thus the divisor range in \(C_m(j)\) contains the
terminal value \(u=N_m(j)=mj+h\).  Its coefficient is
\begin{equation}
 a_R(N_m(j))=-\mu(N_m(j))\log N_m(j),
 \label{eq:tpc34-intro-terminal}
\end{equation}
and vanishes when \(\mu(N_m(j))=0\).  Together with the row sign \(\mu(d)\), this produces
\(\mu(d)\mu(\ell dj+h)\).  The same-time energy consequently contains an
average of four M\"obius factors after expansion.  Existing fixed-form or
logarithmically averaged Chowla results do not directly provide the
growing-coefficient, dyadic, physically masked estimate
\eqref{eq:tpc34-intro-V-gate}; see the precise boundary already recorded
in TPC-33 and the primary sources cited there
\citep{Tao2016LogChowla,MatomakiRadziwillTao2015,WangTPC33}.

An abstract sign-blind model satisfying the same cardinality and
pointwise envelopes has \(E\asymp Q^3J^2\).  It follows that cardinality,
equal-pair deletion, and pointwise envelopes alone cannot establish
\eqref{eq:tpc34-inherited-Q3-gate} uniformly over that abstract class.
This statement is deliberately methodological.  It is not a physical
lower bound, and it leaves open a proof using the fused M\"obius signs or
cancellation between terminal and proper-divisor layers, as well as a
proof using additional arithmetic incidence.  The appearance of this
sign-sensitive boundary is consistent with the classical parity problem
in sieve theory \citep{HeathBrown1982Parity}.

\subsection{Organization and claim boundary}

\Cref{sec:tpc34-physical-interface} imports the physical definitions.
\Cref{sec:tpc34-exact-gram} proves the exact Gram expansion and its
divisor-incidence kernel bound.  \Cref{sec:tpc34-unconditional-blocks}
closes the Gram-copy diagonal and target collision, proves the shared-ultra and
two-time content estimates, and distinguishes savings from closure.
\Cref{sec:tpc34-orbit-sliced} proves the orbit-sliced transfer.
\Cref{sec:tpc34-terminal-barrier} isolates the terminal M\"obius layer and
the sign-blind obstruction.  \Cref{sec:tpc34-high-beta} gives the exact
endpoint ledger and the next theorem menu.  The accompanying certificate
checks only finite algebra and rational bookkeeping.

No estimate of the off-diagonal energy in
\eqref{eq:tpc34-intro-V-gate} is proved.  In particular, this paper does
not prove an unweighted affine Chowla theorem, the zero-frequency
flatness statement of TPC-32, a Hardy--Littlewood prime-pair asymptotic,
positivity, infinitely many twin primes, or a breach of sieve parity.
